\documentclass[12pt]{article}

\input{../../_assets/preambulo.tex}


\begin{document}

    % 1. Foto de fondo
    % 2. Título
    % 3. Encabezado Izquierdo
    % 4. Color de fondo
    % 5. Coord x del titulo
    % 6. Coord y del titulo
    % 7. Fecha

    
    \input{../../_assets/portada}
    \portadaExamen{ffccA4.jpg}{Curvas y\\Superficies\\Examen III}{Curvas y Superficies. Examen III}{MidnightBlue}{-8}{28}{2026}{José Juan Urrutia Milán}

    \begin{description}
        \item[Asignatura] Curvas y Superficies.
        \item[Curso Académico] 2022/23.
        \item[Grado] Grado en Matemáticas.
        % \item[Grupo] B.
        % \item[Profesor] ---.
        \item[Descripción] Convocatoria Ordinaria.
        % \item[Fecha] 15 de junio de 2022.
        % \item[Duración] 2 horas y media.
    \end{description}
    \newpage


    % ------------------------------------

    \noindent
    Pueden faltar algunos ejercicios del examen original.
    
    \begin{ejercicio}
        Sea $\alpha:I\to\bb{S}^2$ una curva p.p.a. y embebida donde $I\subset\mathbb{R}$ es un intervalo abierto. Se llama cono sobre $\alpha$ a $S_\alpha = X(I\times \mathbb{R}^+)$ donde $X:I\times\mathbb{R}^+\to\mathbb{R}^3$ es la aplicación dada por $X(u,v) = v\alpha(u)$.
        \begin{enumerate}[label=\alph*)]
            \item Demuestra que $S_\alpha$ es una superficie regular orientable y determina la primera forma fundamental de $S$ en cualquier punto $p\in S_\alpha$.
            \item Explicita una aplicación de Gauss de $S_\alpha$, y determina la segunda forma fundamental y el endomorfismo de Weingarten de $S_\alpha$ asociados a $N$. Prueba que $S_\alpha$ es llana y calcula la curvatura media para todo $p\in S_\alpha$.
            \item Calcula las curvaturas y direcciones principales para todo $p\in S_\alpha$.
            \item Prueba que la aplicación $\Phi:S_\alpha\to \{z=0\}$ dada por $\Phi(X(u,v)) = v(\cos(u),\sen(u),0)$ es una isometría local. Sea $\beta:I\to\bb{S}^2$ otra curva p.p.a. y embebida. Demuestra que los conos $S_\alpha$ y $S_\beta$ son superficies isométricas.
        \end{enumerate}
    \end{ejercicio}

    \begin{ejercicio}
        Realice los siguientes apartados:
        \begin{enumerate}[label=\alph*)]
            \item Sea $X:\mathbb{R}^+\times\mathbb{R}^+\to S$ una parametrización de una superficie regular $S$ tal que la primera forma fundamental viene dada por:
                \begin{equation*}
                    E(u,v) = u^2, \quad F(u,v) =0,\quad G(u,v) = v^2
                \end{equation*}
                Demuestra que $X(U)$ es isométrica a un abierto del plano.
            \item Sea $Y:\left]-\frac{\pi}{2},\frac{\pi}{2}\right[\times \mathbb{R}\to S$ una parametrización de uan superficie regular $S$ tal que la primera y segunda formas fundamentales vienen dadas por:
                \begin{gather*}
                    E(u,v) = c\sen(u)^2 + \cos(u)^2, \quad F(u,v) = 0, \quad G(u,v) = c\cos(u)^2 \\
                    e(u,v) = \frac{1}{h(u,v)}, \quad f(u,v) = 0, \quad g(u,v) = \frac{\cos(u)^2}{h(u,v)}
                \end{gather*}
                donde $c$ es una constante positiva y $h:\left]-\frac{\pi}{2},\frac{\pi}{2}\right[\times\mathbb{R}\to\mathbb{R}$ es una función positiva. Determina para qué constantes $c$ y funciones $h$ la superficie $Y\left(\left]-\frac{\pi}{2},\frac{\pi}{2}\right[\right)$ está contenido en alguna esfera.
        \end{enumerate}
    \end{ejercicio}

\end{document}
